\itemhead{p.42 练习 · 第2题（教材p.42）}{XJ-B1-C02-PR-P042-Y1011}
\label{XJ-B1-C02-PR-P042-Y1011}
{\footnotesize 湘教版教材p.42 / PDF第49页 / p.42 练习 / 第2题}
\field{原题}{如图，动物园要以墙体为背面，用钢筋网围成四间面积相同的矩形虎笼。
\bankfig{C02-PR042}{0.58}
\subq{1}现有可围 $36\,\mathrm m$ 长钢筋网的材料，每间虎笼的长、宽各设计为多少时，可使每间虎笼的面积最大？
\subq{2}若每间虎笼的面积为 $20\,\mathrm{m^2}$，则每间虎笼的长、宽各设计为多少时，可使围成四间虎笼的钢筋网总长最小？}
\field{参考答案}{设每间沿墙方向宽 $w$、垂直墙长 $l$，钢筋网长 $4w+5l$。（1）$4w+5l=36$ 时 $l=3.6,w=4.5$ 米，每间最大面积16.2平方米；（2）$lw=20$ 时 $l=4,w=5$ 米，总长最小40米。两小问与答案册放大原页一致；方向名称按图中沿墙、垂直墙明确，避免混淆长宽。}
{\small \field{关联知识}{I03 基本不等式、等号与最值建模（教材帮p.52、53、54 / PDF61、62、63）；F03 函数求值、解析式与函数建模（教材帮p.80、87、88、89、90、131 / PDF89、96、97、98、99、140）}
\field{方法与易错点}{列正数条件与固定量，构造基本不等式，解等号条件并代回实际限制。 易错：把一个下界误当最小值；多个等号条件不能同时实现；漏单位和整数限制。}
\field{答案依据}{答案册相关题组：PDF第11页（册内页10）；按当前题库原题逐项重算，不按旧版题号机械对套。}}
